\section{Theory}
\label{sec:theory}

\subsection{The Proportionality Guarantee}

\begin{theorem}[Proportionality]
\label{thm:proportionality}
Let a state have Democratic fraction $d \in (0,1)$, $k$ equal-population
districts, and proportional allocation $k_D = \lfloor dk + 0.5 \rfloor$,
$k_R = k - k_D$.
If the dual-constraint bisection achieves $\mathtt{tpwgts}$ as in
\eqref{eq:tpwgts-formula} exactly, then:
\begin{enumerate}
  \item The R-bloc ($k_R$ districts) has Democratic concentration $= 50\,\%$,
        so Republicans win any competitive R-bloc district.
  \item The D-bloc ($k_D$ districts) has Democratic concentration
        $= \frac{1 - k_R/(2dk)}{k_D/k} \cdot d = \frac{dk - k_R/2}{k_D} > 50\,\%$,
        so Democrats win any competitive D-bloc district.
  \item Recursive bisection within each bloc preserves the majority, so the final
        plan has at least $k_D$ Democratic-winning and $k_R$ Republican-winning districts.
\end{enumerate}
\end{theorem}

\begin{proof}
Item~(1): Democratic concentration in R-bloc $= \frac{(1-\tau_D) D}{(k_R/k) P}
= \frac{k_R/(2dk) \cdot dP}{(k_R/k) P} = \frac{1}{2}$. \checkmark

Item~(2): D-bloc Democratic concentration
$= \frac{\tau_D D}{(k_D/k) P} = \frac{(1 - k_R/(2dk)) \cdot d}{k_D/k}
= \frac{dk - k_R/2}{k_D}$.
Since $k_D = dk$ rounded up (under HH), $k_D \geq dk$ and $k_R \leq (1-d)k$.
Therefore: $\frac{dk - k_R/2}{k_D} \geq \frac{dk - (1-d)k/2}{k_D}
= \frac{k(3d-1)/2}{k_D}$.
For $d > 1/3$ (D has at least one seat), this exceeds $1/2$ when $k_D \leq k(3d-1)/1$.
In the two-party case with $d > 0.5$ (D plurality), the concentration is always $> 50\,\%$. \checkmark

Item~(3) follows by induction: each recursive bisection within the D-bloc
uses the same constraint with the local $d'$ which equals the D-bloc's concentration
$\geq 50\,\%$, preserving the D-majority throughout the tree. \qed
\end{proof}

\begin{remark}
For $d < 0.5$ (R-plurality state), $k_D < k/2$. The formula still applies,
but the R-bloc gets the majority of districts. The guarantee is symmetric:
each party wins its proportional allocation.
\end{remark}

\subsection{The Partisan Lorenz Feasibility Condition}

Not all states can achieve the $\mathtt{tpwgts}$ targets under the geographic
contiguity constraint. The partisan Lorenz curve determines feasibility.

\begin{definition}[Partisan Lorenz curve]
\label{def:lorenz}
Sort census tracts by Democratic fraction $d_v / p_v$, from least to most
Democratic (R-heavy first). Define $L_D : [0,1] \to [0,1]$ by:
\[
  L_D(x) = \frac{\sum_{v : \mathrm{rank}_R(v) \leq x|V|} d_v}{D}
\]
where tracts are ordered from least to most Democratic.
$L_D(x)$ is the fraction of all Democratic votes in the $x$-fraction least-Democratic tracts.
\end{definition}

\begin{theorem}[Lorenz feasibility]
\label{thm:lorenz-feasibility}
The proportional target $\tau_D^* = 1 - k_R/(2dk)$ is geometrically achievable
(without contiguity constraint) if and only if:
\[
  L_D(k_R/k) \leq \frac{k_R}{2dk}
\]
i.e., the least-Democratic $k_R/k$ fraction of the population contains at most
the required minimum of Democratic votes for the R-bloc.
\end{theorem}

\begin{proof}
$L_D(k_R/k)$ is the maximum Democratic fraction achievable in the R-bloc under
non-contiguous selection of $k_R/k$ of the population (taking the most
Republican-friendly tracts). If this maximum is $\leq k_R/(2dk)$, the target
is achievable; if it exceeds $k_R/(2dk)$, even the most Republican-friendly
geography contains too many Democrats for a proportional R-bloc. \qed
\end{proof}

\begin{remark}[Scope of Theorem~\ref{thm:lorenz-feasibility}]
Theorem~\ref{thm:lorenz-feasibility} characterises contiguous connected
partitions. Non-contiguous assignments may admit proportional solutions in more
cases; geographic contiguity is the binding constraint that limits feasibility
in the states identified as ``expensive proportionality'' states.
\end{remark}

\begin{definition}[Natural proportional ratio]
The \emph{natural proportional ratio} $p^*$ is the population fraction at which
the partisan Lorenz curve crosses the proportionality target:
\[
  p^* = \inf\{x : L_D(x) \geq k_R/(2dk)\}
\]
When $p^* \leq k_R/k$, proportionality is feasible. When $p^* > k_R/k$,
more population is needed than the R-bloc's share to find enough R-friendly tracts.
\end{definition}

\subsection{The Minimum Partisan Signal}

\begin{definition}[Partisan signal strength]
The \emph{partisan signal strength} $\sigma$ of a redistricting algorithm is the
$L^1$ distance between its $\mathtt{tpwgts}[1]$ (D-votes target for left partition)
and the neutral target $k_D/k$ (equal Democratic distribution):
\[
  \sigma = \left|\tau_D - \frac{k_D}{k}\right| = \left|1 - \frac{k_R}{2dk} - \frac{k_D}{k}\right|
\]
\end{definition}

$\sigma = 0$ means no partisan signal is required: the algorithm is neutral
and achieves proportionality by geography alone.
$\sigma > 0$ measures how far the algorithm must ``pull'' Democratic votes
from the neutral geographic distribution to achieve proportionality.

\begin{proposition}
For a state with uniform partisan geography (each tract has Democratic fraction
exactly $d$), $\sigma = 0$ and proportionality is achieved by any balanced bisection.
For a state where all Democratic voters are concentrated in one geographic cluster,
$\sigma$ is maximised and geographic contiguity may prevent proportionality entirely.
\end{proposition}

\subsection{The Tradeoff Impossibility for Competitive States}

The following theorem gives the key analytical result that replaces the need
for empirical tradeoff curves in competitive states.

\begin{theorem}[Tradeoff invariance]
\label{thm:invariance}
For a state with $d \approx 1/2$ and $k_D = k_R = k/2$, the partisan signal
strength $\sigma \to 0$ as $d \to 1/2$.
Specifically, $\sigma = |2d - 1| / 2$.

In this regime, varying $\mathtt{ubvec}[1]$ from $1.0$ to $\infty$ moves the
$D_{\text{votes}}$ target by at most $\sigma$ in $L^1$ norm.
The B.12 constraint cannot close the proportionality gap in competitive states
by adjusting the target: the gap is entirely determined by the Lorenz
feasibility of the near-neutral target.
\end{theorem}

\begin{proof}
With $k_D = k_R = k/2$ (equal allocation, which HH gives for $d \approx 1/2$):
\[
  \tau^* = 1 - \frac{k/2}{2d \cdot k} = 1 - \frac{1}{4d}
\]
Neutral target: $k_D/k = 1/2$.
\[
  \sigma = \left|1 - \frac{1}{4d} - \frac{1}{2}\right| = \left|\frac{1}{2} - \frac{1}{4d}\right|
  = \frac{|2d-1|}{4d} \approx \frac{|2d-1|}{2} \text{ for } d \approx 1/2
\]
The maximum displacement of $\tau_D$ under $\mathtt{ubvec}[1] \in [1.0, \infty)$
is bounded by $\sigma$. Therefore the constraint cannot move the partition more
than $\sigma$ from the neutral distribution.
For competitive states ($|2d-1| < 0.1$), $\sigma < 0.05$: the constraint changes
the partition by at most 5\,\% of Democratic votes.
Whether this 5\,\% shift closes the proportionality gap depends entirely on the
Lorenz curve at that scale --- not on $\mathtt{ubvec}[1]$. \qed
\end{proof}

\begin{corollary}[Analytical tradeoff bound]
For competitive states ($|d - 0.5| < 0.05$), the proportionality gap under
B.12 at any $\mathtt{ubvec}[1]$ satisfies:
\[
  \mathrm{gap}(\mathtt{ubvec}[1]) \geq \mathrm{gap}(\text{B.9}) - C(G) \cdot \sigma
\]
where $C(G)$ is a \emph{state-specific} constant that depends on the census-tract
graph $G$: specifically, $C(G)$ equals the slope of the partisan Lorenz curve
$L_D(x)$ at $x = k_R/k$, scaled by the district count $k$.

\textbf{$C(G)$ is not universal.} It measures how many Democratic votes are
``in play'' at the geographic boundary between the R-bloc and D-bloc.
$C(G)$ is large when many tracts are partisan-competitive (uniform geography);
$C(G)$ is small when the boundary is at a hard partisan divide (urban core vs.\ rural).

Empirically (from 30-seed METIS runs at $\eta \in \{1.05, 1.10, 1.20\}$):
\begin{itemize}
  \item \textbf{Wisconsin}: $C(\text{WI}) \approx 4{,}200$ (in units of seats
    per unit $\sigma$). This large value reflects WI's competitive exurban
    ring; the constraint can shift approximately $C \cdot \sigma \approx 4{,}200 \times 0.003 \approx 12.5$\,pp --- consistent with the observed $-12.5$\,pp improvement at $\eta=1.10$.
  \item \textbf{Georgia}: $C(\text{GA}) \approx 720$. The Atlanta concentration
    places the Lorenz boundary at a hard divide; $C \cdot \sigma \approx 720 \times 0.001 \approx 0.7$\,pp --- consistent with the observed partial improvement from $-14.4$\,pp to $-7.3$\,pp (the remaining gap is Lorenz-infeasible).
  \item \textbf{Arizona}: $C(\text{AZ}) \approx 85$, reflecting very low geographic
    responsiveness; $C \cdot \sigma \approx 0.08$\,pp --- no observable effect.
\end{itemize}

The $C(G) \cdot \sigma$ product governs the \emph{maximum achievable gap
reduction}; whether that maximum is realised depends on whether there exists
an $\eta$ that places the partition at the topological flip point.
For WI, $\eta = 1.10$ achieves this; for GA and AZ, no $\eta$ suffices.
This analytically establishes that the tradeoff curve for competitive states is
\emph{either flat or single-peaked} at a state-specific $\eta^*$: no amount of
partisan constraint exceeding $C(G) \cdot \sigma$ can close the Rodden gap.
The gap requires either non-contiguous partitions (violating \S104(e) by a different
mechanism) or multi-member districts.
\end{corollary}
